% GENERATED FILE. Edit canonical lessons, then run npm run generate:books.
% canonical-source-hash: fnv1a64:1c233dae889a5b02
% canonical-lessons: TA-C304-pataru, TA-C304-pitunku, TA-C304-takar, TA-C304-ulukku, TA-C304-citai

\chapter{Shaken and Shattered}
\label{ch:ta-shaken-and-shattered}
\hlchaptermodality{304}

\begin{chapteropening}
\textbf{By the end of this chapter:} \emph{I can say how people and things were shaken.}

Complete the last lesson of chapter 304: I can say how people and things were shaken.
\end{chapteropening}

\section[\texorpdfstring{\ta{பதறு} (pataṟu)}{பதறு (pataṟu)}]{\ta{பதறு} (pataṟu) — to get flustered}

\label{lesson:TA-C304-pataru}



\begin{quote}
Before the new one: say the Tamil for refuge, then the Tamil for to drive away.
\end{quote}

\subsection*{You'll want to know: \ta{பதறு}}
\textbf{\ta{பதறு}} — \emph{pataṟu} — \textquotedblleft{}to get flustered\textquotedblright{}.

\begin{quote}
\ta{சத்தம்} \ta{கேட்டதும்} \ta{எல்லோரும்} \ta{பதறினார்கள்}.
\end{quote}

\begin{quote}
\emph{cattam kēṭṭatum ellōrum pataṟiṉārkaḷ.}
\end{quote}

\begin{quote}
Everyone got flustered when they heard the noise.
\end{quote}

Past (I …): \ta{பதறினேன்}, \emph{pataṟiṉēṉ}.

\ta{பதறாதே}, do not panic: a common way to calm someone.

\begin{cousinweb}
\textbf{\ta{பதறு}} (\emph{pataṟu}) is native Dravidian (DEDR 3910): to be flurried, confused, overhasty.
\end{cousinweb}

\begin{grammarlens}[title={What you've built}]
One more word for telling what happened.
\end{grammarlens}

\subsection*{Your turn}
\begin{itemize}
\raggedright
  \item \emph{Say it:} \emph{pataṟu}
  \item \emph{Say it:} \emph{pataṟu}, once more
  \item \emph{Say it:} \emph{cattam kēṭṭatum ellōrum pataṟiṉārkaḷ.} — Everyone got flustered when they heard the noise.
  \item \emph{Recall:} say the Tamil for refuge, then the Tamil for to drive away, then say \emph{pataṟu} again
\end{itemize}

\begin{tcolorbox}[breakable,colback=teal!4,colframe=teal!35!black,title={Before you move on}]
What does \ta{பதறு} mean? (\textquotedblleft{}To get flustered\textquotedblright{}.) Say it once more.
\end{tcolorbox}
\section[\texorpdfstring{\ta{பிடுங்கு} (piṭuṅku)}{பிடுங்கு (piṭuṅku)}]{\ta{பிடுங்கு} (piṭuṅku) — to snatch; to pull out}

\label{lesson:TA-C304-pitunku}



\begin{quote}
Before the new one: say the Tamil for to drive away, then the Tamil for to get flustered.
\end{quote}

\subsection*{You'll want to know: \ta{பிடுங்கு}}
\textbf{\ta{பிடுங்கு}} — \emph{piṭuṅku} — \textquotedblleft{}to snatch; to pull out\textquotedblright{}.

\begin{quote}
\ta{ஒருவன்} \ta{அவளுடைய} \ta{கைப்பையைப்} \ta{பிடுங்கிக்கொண்டு} \ta{ஓடினான்}.
\end{quote}

\begin{quote}
\emph{oruvaṉ avaḷuṭaiya kaippaiyaip piṭuṅkikkoṇṭu ōṭiṉāṉ.}
\end{quote}

\begin{quote}
A man snatched her handbag and ran off.
\end{quote}

Past (I …): \ta{பிடுங்கினேன்}, \emph{piṭuṅkiṉēṉ}.

Also to pull out a weed or a tooth.

\begin{cousinweb}
\textbf{\ta{பிடுங்கு}} (\emph{piṭuṅku}) is native Dravidian (DEDR 4152): to pull out, pluck up.
\end{cousinweb}

\begin{grammarlens}[title={What you've built}]
One more word for telling what happened.
\end{grammarlens}

\subsection*{Your turn}
\begin{itemize}
\raggedright
  \item \emph{Say it:} \emph{piṭuṅku}
  \item \emph{Say it:} \emph{piṭuṅku}, once more
  \item \emph{Say it:} \emph{oruvaṉ avaḷuṭaiya kaippaiyaip piṭuṅkikkoṇṭu ōṭiṉāṉ.} — A man snatched her handbag and ran off.
  \item \emph{Recall:} say the Tamil for to drive away, then the Tamil for to get flustered, then say \emph{piṭuṅku} again
\end{itemize}

\begin{tcolorbox}[breakable,colback=teal!4,colframe=teal!35!black,title={Before you move on}]
What does \ta{பிடுங்கு} mean? (\textquotedblleft{}To snatch; to pull out\textquotedblright{}.) Say it once more.
\end{tcolorbox}
\section[\texorpdfstring{\ta{தகர்} (takar)}{தகர் (takar)}]{\ta{தகர்} (takar) — to smash; to shatter}

\label{lesson:TA-C304-takar}



\begin{quote}
Before the new one: say the Tamil for to get flustered, then the Tamil for to snatch.
\end{quote}

\subsection*{You'll want to know: \ta{தகர்}}
\textbf{\ta{தகர்}} — \emph{takar} — \textquotedblleft{}to smash; to shatter\textquotedblright{}.

\begin{quote}
\ta{திருடர்கள்} \ta{கதவின்} \ta{பூட்டைத்} \ta{தகர்த்தார்கள்}.
\end{quote}

\begin{quote}
\emph{tiruṭarkaḷ kataviṉ pūṭṭait takarttārkaḷ.}
\end{quote}

\begin{quote}
The thieves smashed the lock of the door.
\end{quote}

Past (I …): \ta{தகர்த்தேன்}, \emph{takarttēṉ}.

Without an object: \ta{தகர்ந்தது}, it was shattered.

\begin{cousinweb}
\textbf{\ta{தகர்}} (\emph{takar}) is native Dravidian (DEDR 2999): to be broken to pieces, and with -\ta{த்த}- to smash.
\end{cousinweb}

\begin{grammarlens}[title={What you've built}]
One more word for telling what happened.
\end{grammarlens}

\subsection*{Your turn}
\begin{itemize}
\raggedright
  \item \emph{Say it:} \emph{takar}
  \item \emph{Say it:} \emph{takar}, once more
  \item \emph{Say it:} \emph{tiruṭarkaḷ kataviṉ pūṭṭait takarttārkaḷ.} — The thieves smashed the lock of the door.
  \item \emph{Recall:} say the Tamil for to get flustered, then the Tamil for to snatch, then say \emph{takar} again
\end{itemize}

\begin{tcolorbox}[breakable,colback=teal!4,colframe=teal!35!black,title={Before you move on}]
What does \ta{தகர்} mean? (\textquotedblleft{}To smash; to shatter\textquotedblright{}.) Say it once more.
\end{tcolorbox}
\section[\texorpdfstring{\ta{உலுக்கு} (ulukku)}{உலுக்கு (ulukku)}]{\ta{உலுக்கு} (ulukku) — to shake up}

\label{lesson:TA-C304-ulukku}



\begin{quote}
Before the new one: say the Tamil for to snatch, then the Tamil for to smash.
\end{quote}

\subsection*{You'll want to know: \ta{உலுக்கு}}
\textbf{\ta{உலுக்கு}} — \emph{ulukku} — \textquotedblleft{}to shake up\textquotedblright{}.

\begin{quote}
\ta{அந்தச்} \ta{செய்தி} \ta{ஊர்} \ta{முழுவதையும்} \ta{உலுக்கியது}.
\end{quote}

\begin{quote}
\emph{antac ceyti ūr muḻuvataiyum ulukkiyatu.}
\end{quote}

\begin{quote}
That news shook up the whole town.
\end{quote}

Past (I …): \ta{உலுக்கினேன்}, \emph{ulukkiṉēṉ}.

Also literally: \ta{மரத்தை} \ta{உலுக்கு}, to shake a tree.

\begin{cousinweb}
\textbf{\ta{உலுக்கு}} (\emph{ulukku}) is native Dravidian (DEDR 1003): to shake, as a tree, or to tremble, as in an earthquake.
\end{cousinweb}

\begin{grammarlens}[title={What you've built}]
One more word for telling what happened.
\end{grammarlens}

\subsection*{Your turn}
\begin{itemize}
\raggedright
  \item \emph{Say it:} \emph{ulukku}
  \item \emph{Say it:} \emph{ulukku}, once more
  \item \emph{Say it:} \emph{antac ceyti ūr muḻuvataiyum ulukkiyatu.} — That news shook up the whole town.
  \item \emph{Recall:} say the Tamil for to snatch, then the Tamil for to smash, then say \emph{ulukku} again
\end{itemize}

\begin{tcolorbox}[breakable,colback=teal!4,colframe=teal!35!black,title={Before you move on}]
What does \ta{உலுக்கு} mean? (\textquotedblleft{}To shake up\textquotedblright{}.) Say it once more.
\end{tcolorbox}
\section[\texorpdfstring{\ta{சிதை} (citai)}{சிதை (citai)}]{\ta{சிதை} (citai) — to ruin; to shatter}

\label{lesson:TA-C304-citai}



\begin{quote}
Before the new one: say the Tamil for to smash, then the Tamil for to shake up.
\end{quote}

\subsection*{You'll want to know: \ta{சிதை}}
\textbf{\ta{சிதை}} — \emph{citai} — \textquotedblleft{}to ruin; to shatter\textquotedblright{}.

\begin{quote}
\ta{வெள்ளம்} \ta{விவசாயிகளின்} \ta{கனவுகளைச்} \ta{சிதைத்தது}.
\end{quote}

\begin{quote}
\emph{veḷḷam vivacāyikaḷiṉ kaṉavukaḷaic citaittatu.}
\end{quote}

\begin{quote}
The flood ruined the farmers' dreams.
\end{quote}

Past (I …): \ta{சிதைத்தேன்}, \emph{citaittēṉ}.

Without an object: \ta{சிதைந்தது}, it was ruined or broke apart.

\begin{cousinweb}
\textbf{\ta{சிதை}} (\emph{citai}) is native Dravidian (DEDR 2526): to be spoiled or broken, and with -\ta{த்த}- to spoil, to ruin.
\end{cousinweb}

\begin{grammarlens}[title={What you've built}]
One more word for telling what happened.
\end{grammarlens}

\subsection*{Your turn}
\begin{itemize}
\raggedright
  \item \emph{Say it:} \emph{citai}
  \item \emph{Say it:} \emph{citai}, once more
  \item \emph{Say it:} \emph{veḷḷam vivacāyikaḷiṉ kaṉavukaḷaic citaittatu.} — The flood ruined the farmers' dreams.
  \item \emph{Recall:} say the Tamil for to smash, then the Tamil for to shake up, then say \emph{citai} again
\end{itemize}

\begin{tcolorbox}[breakable,colback=teal!4,colframe=teal!35!black,title={Before you move on}]
What does \ta{சிதை} mean? (\textquotedblleft{}To ruin; to shatter\textquotedblright{}.) Say it once more.
\end{tcolorbox}
